\chapter{Évaluation du Module 2}
\label{chap:evalmod2}

Cette évaluation valide les compétences liées à la méthodologie d'audit, la compréhension des référentiels (ISO 19011, 27001, 9001) et la capacité à rédiger un rapport professionnel.

% ------------------------------------------------------------------------------
\section{Partie A : QCM (25 Questions)}
% ------------------------------------------------------------------------------

\textbf{Consigne :} Cochez la ou les bonnes réponses.

\begin{enumerate}
    \item L'indépendance de l'auditeur signifie :
    \begin{itemize}
        \item[A.] Qu'il ne doit pas être ami avec l'audité.
        \item[B.] Qu'il ne doit pas être juge et partie (pas d'intérêt direct dans le résultat).
        \item[C.] Qu'il doit travailler seul.
        \item[D.] Qu'il doit être externe à l'entreprise.
    \end{itemize}

    \item La norme qui donne les lignes directrices pour réaliser un audit est :
    \begin{itemize}
        \item[A.] ISO 27001.
        \item[B.] ISO 19011.
        \item[C.] ISO 9001.
        \item[D.] ISO 27005.
    \end{itemize}

    \item Une "Non-Conformité Majeure" est :
    \begin{itemize}
        \item[A.] Un petit incident isolé.
        \item[B.] Une absence totale de système ou une défaillance généralisée.
        \item[C.] Une suggestion d'amélioration.
        \item[D.] Une mesure de sécurité non applicable.
    \end{itemize}

    \item La Déclaration d'Applicabilité (SOA) est un document lié à :
    \begin{itemize}
        \item[A.] L'ISO 9001.
        \item[B.] L'ISO 27001.
        \item[C.] L'ISO 19011.
        \item[D.] COBIT.
    \end{itemize}

    \item La "Correction" consiste à :
    \begin{itemize}
        \item[A.] Éliminer la cause du problème.
        \item[B.] Éliminer le problème immédiat (symptôme).
        \item[C.] Blâmer le responsable.
        \item[D.] Écrire une procédure.
    \end{itemize}

    \item Qui est le "Commanditaire" de l'audit ?
    \begin{itemize}
        \item[A.] La personne auditée.
        \item[B.] Celui qui demande et paie l'audit.
        \item[C.] L'auditeur principal.
        \item[D.] Le certificateur.
    \end{itemize}

    \item La "Réunion d'Ouverture" sert à :
    \begin{itemize}
        \item[A.] Présenter les résultats de l'audit.
        \item[B.] Confirmer le plan d'audit et les règles du jeu.
        \item[C.] Former les employés.
        \item[D.] Signer le certificat ISO.
    \end{itemize}

    \item L'Annexe A de l'ISO 27001 contient :
    \begin{itemize}
        \item[A.] Les exigences du système de management.
        \item[B.] La liste des mesures de sécurité (contrôles).
        \item[C.] Les procédures d'audit.
        \item[D.] La politique de sécurité type.
    \end{itemize}

    \item COBIT est un référentiel pour :
    \begin{itemize}
        \item[A.] La gouvernance et le management des SI.
        \item[B.] La qualité des logiciels uniquement.
        \item[C.] La gestion de projet.
        \item[D.] La sécurité physique.
    \end{itemize}

    \item Un constat d'audit doit contenir :
    \begin{itemize}
        \item[A.] Une opinion personnelle.
        \item[B.] Des faits, des preuves et une référence au critère.
        \item[C.] Une proposition de solution immédiate.
        \item[D.] Le nom du coupable.
    \end{itemize}

    \item La High Level Structure (HLS) permet :
    \begin{itemize}
        \item[A.] D'avoir une structure commune pour toutes les normes ISO de management.
        \item[B.] De mieux structurer le code informatique.
        \item[C.] De gérer la sécurité réseau.
        \item[D.] De hiérarchiser les employés.
    \end{itemize}

    \item L'échantillonnage en audit consiste à :
    \begin{itemize}
        \item[A.] Tout vérifier.
        \item[B.] Sélectionner une partie représentative pour tirer des conclusions.
        \item[C.] Choisir les meilleurs éléments.
        \item[D.] Ignorer les anomalies.
    \end{itemize}

    \item L'audit de certification est réalisé par :
    \begin{itemize}
        \item[A.] L'entreprise elle-même.
        \item[B.] Un organisme certificateur externe.
        \item[C.] Un client.
        \item[D.] Le syndicat.
    \end{itemize}

    \item Une "Action Corrective" vise à :
    \begin{itemize}
        \item[A.] Corriger l'erreur immédiatement.
        \item[B.] Empêcher l'erreur de se reproduire (traiter la cause).
        \item[C.] Punir l'employé.
        \item[D.] Transférer le risque.
    \end{itemize}

    \item ISO 9001 traite de :
    \begin{itemize}
        \item[A.] Sécurité de l'information.
        \item[B.] Management de la Qualité.
        \item[C.] Continuité d'activité.
        \item[D.] Protection des données personnelles.
    \end{itemize}

    \item La "Piste d'audit" (Audit Trail) est :
    \begin{itemize}
        \item[A.] Le chemin parcouru par l'auditeur dans les locaux.
        \item[B.] L'enchaînement de preuves permettant de reconstituer un événement.
        \item[C.] Le plan d'audit.
        \item[D.] La facture de l'audit.
    \end{itemize}

    \item Lors d'un entretien d'audit, l'auditeur doit :
    \begin{itemize}
        \item[A.] Poser des questions fermées (Oui/Non) uniquement.
        \item[B.] Poser des questions ouvertes (Comment, Pourquoi) et écouter.
        \item[C.] Donner son avis.
        \item[D.] Intimider l'audité pour obtenir des aveux.
    \end{itemize}

    \item Le rapport d'audit doit être remis :
    \begin{itemize}
        \item[A.] À l'audité uniquement.
        \item[B.] Au commanditaire.
        \item[C.] À la concurrence.
        \item[D.] Au gouvernement.
    \end{itemize}

    \item Un système de management est considéré "Mature" (Niveau 5 CMMI) si :
    \begin{itemize}
        \item[A.] Il n'existe pas.
        \item[B.] Il est en amélioration continue optimisée.
        \item[C.] Il est documenté.
        \item[D.] Il est géré par intuition.
    \end{itemize}

    \item Le PDCA (Plan-Do-Check-Act) est :
    \begin{itemize}
        \item[A.] Une méthode de hacking.
        \item[B.] Un cycle d'amélioration continue.
        \item[C.] Un type de pare-feu.
        \item[D.] Un logiciel de compta.
    \end{itemize}

    \item Si l'auditeur n'a pas accès aux documents demandés, c'est :
    \begin{itemize}
        \item[A.] Une observation.
        \item[B.] Une opportunité.
        \item[C.] Une non-conformité (entrave à l'audit).
        \item[D.] Un droit de l'audité.
    \end{itemize}

    \item La "Réunion de Clôture" se déroule :
    \begin{itemize}
        \item[A.] Avant l'audit.
        \item[B.] Au milieu de l'audit.
        \item[C.] À la fin de la phase terrain, avant le rapport final.
        \item[D.] Après la certification.
    \end{itemize}

    \item La "Checklist" est :
    \begin{itemize}
        \item[A.] La liste des courses de l'auditeur.
        \item[B.] La feuille de route listant les points à vérifier.
        \item[C.] Le rapport final.
        \item[D.] La liste des non-conformités.
    \end{itemize}

    \item Une "Observation" dans un rapport d'audit est :
    \begin{itemize}
        \item[A.] Une non-conformité majeure.
        \item[B.] Une piste d'amélioration qui n'est pas une non-conformité stricte.
        \item[C.] Une faute grave.
        \item[D.] Une description de l'entreprise.
    \end{itemize}

    \item L'audit interne est souvent réalisé pour :
    \begin{itemize}
        \item[A.] Remplacer la DSI.
        \item[B.] Préparer une certification ou vérifier la conformité interne.
        \item[C.] Licencier des employés.
        \item[D.] Espionner les concurrents.
    \end{itemize}
\end{enumerate}

% ------------------------------------------------------------------------------
\section{Partie B : Questions Ouvertes (10 Questions)}
% ------------------------------------------------------------------------------

\textbf{Consigne :} Répondez de manière synthétique (5-10 lignes par question).

\begin{enumerate}
    \item Expliquez la différence fondamentale entre une "Correction" et une "Action Corrective" avec un exemple.
    \item Pourquoi l'indépendance de l'auditeur est-elle une règle absolue ? Que se passe-t-il si elle n'est pas respectée ?
    \item Quelle est l'utilité de la "Déclaration d'Applicabilité" (SOA) dans le cadre d'une certification ISO 27001 ?
    \item Décrivez les 3 sources de preuves possibles lors d'un audit.
    \item Quelle est la différence entre une Non-Conformité Majeure et Mineure ? Donnez un exemple pour chaque cas.
    \item Expliquez le concept de "High Level Structure" (HLS) et son avantage pour les entreprises.
    \item Dans la norme ISO 19011, à quoi sert la réunion d'ouverture ? Citez 3 points abordés.
    \item Pourquoi l'audit est-il considéré comme une étape "Check" (Vérifier) du cycle PDCA ?
    \item Comparez l'approche de l'ISO 27001 (Sécurité) et de l'ISO 9001 (Qualité) : sont-elles compatibles ?
    \item Quel est le rôle du "Commanditaire" de l'audit et pourquoi est-il important de l'identifier ?
\end{enumerate}

% ------------------------------------------------------------------------------
\section{Partie C : Cas Pratique - Mise en Condition (TP)}
% ------------------------------------------------------------------------------

\subsection{Scénario : L'Audit de la Sauvegarde}

\begin{boxlizbiz}
\textbf{Contexte :}
Suite aux risques identifiés sur le Ransomware (Module 1), la direction de LiZbiZ a demandé un audit ciblé sur le processus de sauvegarde. Le PDG veut être sûr que l'entreprise peut repartir en cas d'attaque.

\textbf{Situation découverte :}
Vous avez accès au Lab GNS3 et vous avez interrogé l'Admin Système.
\begin{itemize}
    \item \textbf{Entretien Admin :} "On sauvegarde tout les soirs sur un NAS interne dans le bureau du DSI."
    \item \textbf{Observation :} Le NAS est bien branché, voyant vert allumé.
    \item \textbf{Test Technique (Lab) :} Vous essayez de lister le contenu de la sauvegarde. Le mot de passe du NAS est "admin123". Les sauvegardes des 3 derniers jours sont absentes (erreur disque plein).
    \item \textbf{Document :} La procédure de sauvegarde dit : "Sauvegardes externalisées et cryptées".
\end{itemize}
\end{boxlizbiz}

\subsection{Travaux à Réaliser (En temps limité)}

\textbf{Livrable 1 : Fiche de Non-Conformité (30 min)}
En vous basant sur l'ISO 27001 (Annexe A.12.3 Sauvegarde des informations) :
\begin{itemize}
    \item Rédigez une fiche de constat complète (Critère, Constat, Preuve, Risque).
    \item Qualifiez la Non-Conformité (Majeure ou Mineure) et justifiez.
\end{itemize}

\textbf{Livrable 2 : Rapport d'Audit "Executive Summary" (30 min)}
Rédigez la page de synthèse du rapport destinée au PDG de LiZbiZ.
\begin{itemize}
    \item Donnez un avis global sur le niveau de maturité de la sauvegarde.
    \item Listez les risques principaux liés à la situation trouvée.
\end{itemize}

\textbf{Livrable 3 : Plan d'Action (30 min)}
Proposez le Plan d'Actions Correctives (PAC) pour remédier à la situation.
\begin{itemize}
    \item Distinguez bien les Corrections (immédiat) des Actions Correctives (fond).
    \item Définissez un RACI pour la mise en œuvre.
\end{itemize}

% ==============================================================================
% Fichier : 04_module_incident.tex (Extrait Chapitre 1)
% Description : Introduction à la Gestion des Incidents
% ==============================================================================